\subsubsection{Subproceso 1: Acreditación y Pre-requisitos}

\textbf{Inicio:} El subproceso empieza en los días previos a la matrícula regular, cuando la universidad cierra las actas de notas del semestre anterior y el estudiante realiza el pago de su autoseguro médico y regulariza cualquier deuda pendiente con su Facultad.

\textbf{Incluye:}
\begin{itemize}
    \item Verificación del pago del autoseguro estudiantil para el semestre que inicia.
    \item Comprobación de que el alumno no tenga deudas de libros, materiales o equipos en bibliotecas y laboratorios de su Facultad.
    \item Consolidación de las notas cerradas del ciclo previo y actualización de cambios de plan de estudios o convalidaciones en el sistema.
    \item Cálculo del promedio ponderado de los dos últimos ciclos regulares, asignación del tope de créditos y ubicación del estudiante en su grupo y horario de matrícula.
    \item Generación de la ficha de matrícula individual y activación del estado habilitado en la Intranet DIRCE.
\end{itemize}

\textbf{Fin:} Termina cuando la Intranet DIRCE muestra al estudiante sin bloqueos ni adeudos, con su pago de autoseguro validado, su ficha de matrícula lista y su grupo con fecha y hora asignados para ingresar al Subproceso 2.

\textbf{Fuera del alcance:}
\begin{itemize}
    \item El armado de horarios, apertura de secciones y asignación de docentes por los Departamentos Académicos.
    \item La elección de cursos, secciones y reserva de vacantes en el portal (Subproceso 2).
    \item La revisión de reglas durante la matrícula, el control del Grupo Cero y la atención de solicitudes de rectificación (Subproceso 3.A).
    \item El cierre oficial de la matrícula, la emisión de la boleta visada y el envío de reportes a SUNEDU y MINEDU (Subproceso 3.B).
\end{itemize}
